\documentclass[11pt,a4paper]{article}
\usepackage[utf8]{inputenc}
\usepackage[T1]{fontenc}
\usepackage{lmodern}
\usepackage[left=2.05cm,right=2.05cm,top=1.95cm,bottom=2.05cm]{geometry}
\usepackage{amsmath,amssymb,booktabs,array,tabularx,graphicx,xcolor,fancyhdr,enumitem}
\usepackage{caption,microtype,eso-pic}
\usepackage[hidelinks]{hyperref}
\renewcommand{\figurename}{Figura}
\renewcommand{\tablename}{Tabla}
\definecolor{azul}{HTML}{233D54}
\definecolor{gris}{HTML}{565E64}
\setlength{\parindent}{0pt}
\setlength{\parskip}{0.40em}
\setlength{\tabcolsep}{3.5pt}
\renewcommand{\arraystretch}{1.16}
\setlist[itemize]{leftmargin=1.2em,itemsep=2pt,topsep=3pt}
\captionsetup{font=small,labelfont=bf,skip=4pt}
\pagestyle{fancy}\fancyhf{}
\fancyhead[L]{\small\color{gris}MCDI501 \;|\; Sumativa 3: modelamiento integrado}
\fancyhead[R]{\small\color{gris}Grupo 7}
\fancyfoot[L]{\footnotesize\color{gris}Universidad Andrés Bello \;|\; Magíster en Ciencia de Datos e IA}
\fancyfoot[R]{\footnotesize\thepage\,/\,10}
\setlength{\headheight}{14pt}
\renewcommand{\headrulewidth}{0.3pt}
\newcommand{\titulo}[2]{\vspace*{0.1em}{\Large\bfseries\color{azul}#1}\par{\small\color{gris}#2}\par\vspace{0.35em}}
\newcommand{\subtitulo}[1]{\par\vspace{0.1em}{\bfseries\color{azul}#1}\par}
\newcommand{\figura}[3]{\begin{center}\includegraphics[width=\linewidth,height=#2,keepaspectratio]{../../results/05_sumativa3/figures/#1}\captionof{figure}{#3}\end{center}}
\newcommand{\fuente}{\vspace{-0.3em}{\footnotesize\color{gris}Fuente: elaboración del Grupo 7; datos y cálculos de esta ejecución.}\par}
\newcommand{\AUCFinal}{0,8619}
\newcommand{\RecallFinal}{68,26}
\newcommand{\PrecisionFinal}{57,98}
\newcommand{\AccuracyFinal}{81,79}
\newcommand{\FUnoFinal}{0,6270}
\newcommand{\PRAUCFinal}{0,6872}
\newcommand{\UmbralFinal}{0,6034}
\newcommand{\FalsosNegativos}{3\,035}
\newcommand{\FalsosPositivos}{4\,732}
\newcommand{\VerdaderosPositivos}{6\,528}
\newcommand{\CambioRecall}{28,44}
\newcommand{\CambioAccuracy}{-1,09}
\newcommand{\PesoYes}{2,2303}
\newcommand{\PesoNo}{0,6445}
\newcommand{\MejoraRMSE}{41,76}
\newcommand{\BootstrapValidas}{2\,000}
\newcommand{\FaltantesSunModelo}{47,69}
\newcommand{\MaxVIF}{3,054}
\newcommand{\CookMarcados}{6\,417}
\newcommand{\ModeloFinal}{M2 stepwise AIC}
\newcommand{\VariablesFinales}{\texttt{Humidity3pm}, \texttt{WindGustSpeed}, \texttt{Pressure3pm}, \texttt{Sunshine}, \texttt{log1p\_Rainfall}, \texttt{Cloud3pm}, \texttt{MaxTemp}, \texttt{Humidity9am}}
\newcommand{\RMSESolReg}{2,226}
\newcommand{\RMSESolMed}{3,822}
\newcommand{\RDSol}{0,645}
\newcommand{\TablaImputacion}{\fontsize{8.3}{10.2}\selectfont\begin{tabularx}{\linewidth}{@{}lXrrrr@{}}\toprule
Variable & Tres predictores & Donantes & $R^2$ val. & RMSE reg. & RMSE med.\\\midrule
MinTemp & Temp9am; Evaporation; Pressure3pm & 56\,512 & 0,827 & 2,69 & 6,48\\
MaxTemp & Temp3pm; MinTemp; Evaporation & 56\,017 & 0,973 & 1,16 & 7,06\\
Rainfall & Humidity3pm; Humidity9am; Sunshine & 51\,140 & 0,081 & 8,50 & 9,18\\
Evaporation & MaxTemp; MinTemp; Humidity9am & 56\,479 & 0,417 & 3,27 & 4,33\\
Sunshine & Humidity3pm; Cloud3pm; Cloud9am & 44\,903 & 0,645 & 2,23 & 3,82\\
WindGustSpeed & WindSpeed3pm; WindSpeed9am; Pressure9am & 85\,703 & 0,617 & 8,27 & 13,41\\
WindSpeed9am & WindGustSpeed; WindSpeed3pm; Humidity9am & 91\,836 & 0,404 & 6,69 & 8,74\\
WindSpeed3pm & WindGustSpeed; WindSpeed9am; Pressure9am & 85\,703 & 0,482 & 6,15 & 8,54\\
Humidity9am & Humidity3pm; Evaporation; Sunshine & 48\,718 & 0,543 & 12,60 & 18,65\\
Humidity3pm & Sunshine; Rainfall; Humidity9am & 51\,140 & 0,588 & 13,13 & 20,45\\
Pressure9am & Pressure3pm; MinTemp; WindGustSpeed & 85\,689 & 0,929 & 1,87 & 7,05\\
Pressure3pm & Pressure9am; Temp9am; WindGustSpeed & 85\,666 & 0,932 & 1,84 & 7,08\\
Cloud9am & Sunshine; Cloud3pm; Humidity3pm & 44\,903 & 0,495 & 1,99 & 2,88\\
Cloud3pm & Sunshine; Cloud9am; Humidity3pm & 44\,903 & 0,541 & 1,80 & 2,74\\
Temp9am & MinTemp; MaxTemp; Evaporation & 56\,732 & 0,928 & 1,78 & 6,65\\
Temp3pm & MaxTemp; MinTemp; Evaporation & 56\,017 & 0,971 & 1,17 & 6,93\\
\bottomrule\end{tabularx}}
\newcommand{\TablaSeleccion}{\small\begin{tabularx}{\linewidth}{@{}lrrrrr@{}}\toprule
Modelo & p & AIC & BIC & AUC OOF & Umbral\\\midrule
M1 & 4 & 80\,289,2 & 80\,336,8 & 0,8213 & 0,5638\\
M2 & 8 & 72\,746,5 & 72\,832,1 & 0,8604 & 0,6034\\
M3 & 3 & 80\,848,3 & 80\,886,3 & 0,8180 & 0,5776\\
\bottomrule\end{tabularx}}
\newcommand{\TablaCoeficientesTres}{\small\begin{tabularx}{\linewidth}{@{}lrrr@{}}\toprule
Predictor & M1 $\beta$ (OR) & M2 $\beta$ (OR) & M3 $\beta$ (OR)\\\midrule
const & -0,415 (0,660) & -0,579 (0,561) & -0,416 (0,659)\\
Humidity3pm & 0,894 (2,446) & 0,902 (2,465) & 0,802 (2,230)\\
Sunshine & -0,660 (0,517) & -0,478 (0,620) & -0,600 (0,549)\\
Rainfall & 0,384 (1,468) & --- & 0,375 (1,455)\\
MaxTemp & 0,272 (1,313) & 0,095 (1,099) & ---\\
WindGustSpeed & --- & 0,467 (1,596) & ---\\
Pressure3pm & --- & -0,486 (0,615) & ---\\
log1p\_Rainfall & --- & 0,258 (1,294) & ---\\
Cloud3pm & --- & 0,193 (1,213) & ---\\
Humidity9am & --- & 0,101 (1,106) & ---\\
\bottomrule\end{tabularx}}
\newcommand{\TablaMetricas}{\small\begin{tabularx}{\linewidth}{@{}llrrrrr@{}}\toprule
Modelo & Muestra & Acc.\% & Prec.\% & Recall\% & F1 & AUC\\\midrule
M1 & train & 77,90 & 50,54 & 66,65 & 0,575 & 0,8214\\
M1 & test & 78,06 & 50,82 & 66,65 & 0,577 & 0,8240\\
M2 & train & 81,55 & 57,52 & 67,65 & 0,622 & 0,8606\\
M2 & test & 81,79 & 57,98 & 68,26 & 0,627 & 0,8619\\
M3 & train & 78,32 & 51,30 & 64,97 & 0,573 & 0,8181\\
M3 & test & 78,57 & 51,76 & 64,82 & 0,576 & 0,8208\\
\bottomrule\end{tabularx}}
\newcommand{\TablaBootstrap}{\fontsize{8.4}{10.3}\selectfont\begin{tabularx}{\linewidth}{@{}lrrrrr@{}}\toprule
Predictor & $\hat\beta$ & IC Wald beta & IC boot beta & OR [IC boot] & Signo\%\\\midrule
Humidity3pm & 0,902 & [0,877; 0,927] & [0,804; 1,025] & 2,465 [2,234; 2,787] & 100,0\\
WindGustSpeed & 0,467 & [0,449; 0,486] & [0,399; 0,547] & 1,596 [1,491; 1,729] & 100,0\\
Pressure3pm & -0,486 & [-0,506; -0,467] & [-0,602; -0,388] & 0,615 [0,548; 0,679] & 100,0\\
Sunshine & -0,478 & [-0,506; -0,451] & [-0,550; -0,401] & 0,620 [0,577; 0,670] & 100,0\\
log1p\_Rainfall & 0,258 & [0,241; 0,275] & [0,194; 0,308] & 1,294 [1,214; 1,360] & 100,0\\
Cloud3pm & 0,193 & [0,169; 0,217] & [0,117; 0,260] & 1,213 [1,124; 1,297] & 100,0\\
MaxTemp & 0,095 & [0,074; 0,116] & [-0,034; 0,212] & 1,099 [0,966; 1,237] & 92,6\\
Humidity9am & 0,101 & [0,077; 0,125] & [0,011; 0,193] & 1,106 [1,011; 1,213] & 98,5\\
\bottomrule\end{tabularx}}
\newcommand{\PredictoresInestables}{\texttt{MaxTemp}}
\newcommand{\NumeroFallosBootstrap}{0}
\newcommand{\LocalidadesInfluyentes}{MountGinini, Wollongong, Ballarat}
\newcommand{\TablaVIF}{\small\begin{tabularx}{\linewidth}{@{}lr@{}}\toprule
Predictor & VIF\\\midrule
Humidity3pm & 2,561\\
WindGustSpeed & 1,325\\
Pressure3pm & 1,522\\
Sunshine & 3,054\\
log1p\_Rainfall & 1,315\\
Cloud3pm & 2,222\\
MaxTemp & 1,856\\
Humidity9am & 2,183\\
\bottomrule\end{tabularx}}
\newcommand{\TablaLinealidad}{\footnotesize\begin{tabularx}{\linewidth}{@{}lrr@{}}\toprule
Predictor & p ajustado & $\Delta$ dev.\%\\\midrule
Humidity3pm & $<0{,}001$ & 0,154\\
WindGustSpeed & 1,000 & 0,001\\
Pressure3pm & 1,000 & 0,001\\
Sunshine & $<0{,}001$ & 0,127\\
log1p\_Rainfall & $<0{,}001$ & 0,132\\
MaxTemp & 0,405 & 0,004\\
Humidity9am & 0,069 & 0,007\\
\bottomrule\end{tabularx}}
\newcommand{\TablaImpacto}{\small\begin{tabularx}{\linewidth}{@{}llrrrrrr@{}}\toprule
Estrategia & Prueba & $n_{train}$ & $n_{test}$ & Loc. & AUC & Acc.\% & Recall\%\\\midrule
Eliminación & Propia & 42\,282 & 18\,234 & 27 & 0,8813 & 82,83 & 70,98\\
Eliminación & Común & 42\,282 & 18\,234 & 27 & 0,8813 & 82,83 & 70,98\\
Mediana & Propia & 99\,535 & 42\,658 & 49 & 0,8589 & 81,68 & 67,75\\
Mediana & Común & 99\,535 & 18\,234 & 27 & 0,8801 & 82,14 & 73,49\\
Regresión & Propia & 99\,535 & 42\,658 & 49 & 0,8619 & 81,79 & 68,26\\
Regresión & Común & 99\,535 & 18\,234 & 27 & 0,8809 & 82,98 & 70,96\\
\bottomrule\end{tabularx}}
\newcommand{\CoberturaEliminacion}{42,74}
\newcommand{\LocEliminacion}{27}
\newcommand{\TablaCoefImpacto}{\footnotesize\begin{tabularx}{\linewidth}{@{}lrrr@{}}\toprule
Predictor & Eliminación & Mediana & Regresión\\\midrule
Humidity3pm & 0,967 [0,926; 1,009] & 1,013 [0,989; 1,036] & 0,902 [0,877; 0,927]\\
WindGustSpeed & 0,519 [0,490; 0,549] & 0,488 [0,470; 0,507] & 0,467 [0,449; 0,486]\\
Pressure3pm & -0,492 [-0,524; -0,461] & -0,489 [-0,508; -0,470] & -0,486 [-0,506; -0,467]\\
Sunshine & -0,488 [-0,520; -0,455] & -0,383 [-0,404; -0,362] & -0,478 [-0,506; -0,451]\\
log1p\_Rainfall & 0,191 [0,164; 0,219] & 0,234 [0,217; 0,251] & 0,258 [0,241; 0,275]\\
Cloud3pm & 0,245 [0,213; 0,276] & 0,287 [0,266; 0,308] & 0,193 [0,169; 0,217]\\
MaxTemp & 0,108 [0,072; 0,143] & 0,067 [0,046; 0,087] & 0,095 [0,074; 0,116]\\
Humidity9am & 0,070 [0,029; 0,111] & 0,148 [0,124; 0,172] & 0,101 [0,077; 0,125]\\
\bottomrule\end{tabularx}}
\newcommand{\TablaMonteCarlo}{\small\begin{tabularx}{\linewidth}{@{}lrrrrr@{}}\toprule
Humedad & $n$ train & p lluvia & Media /30 & Capacidad 95\% & MCSE\\\midrule
Alta & 18\,924 & 0,5834 & 17,52 & 22 & 0,027\\
Menor & 78\,075 & 0,1357 & 4,04 & 7 & 0,019\\
\bottomrule\end{tabularx}}
\newcommand{\CapacidadAlta}{22}
\newcommand{\CapacidadBaja}{7}
\newcommand{\MediaMCAlta}{17,52}
\newcommand{\ProbMCAlta}{58,34}
\newcommand{\MediaMCBaja}{4,04}
\newcommand{\ProbMCBaja}{13,57}
\newcommand{\ORHumedad}{2,465}
\newcommand{\ICHumedad}{[2,234; 2,787]}
\newcommand{\DEHumedad}{20,60}
\newcommand{\ORDoblarRain}{1,219}
\newcommand{\FallbackSol}{35\,960}
\newcommand{\ImputadosSol}{47\,607}

\begin{document}
% Portada editable: estructura y posiciones del PDF de Sumativa 1 del Grupo 7.
% Las coordenadas están en bp, medidas desde la esquina superior izquierda.
% El logotipo es el recurso original extraído sin alteraciones de ese PDF.
\newcommand{\portadaEvaluacion}{Evaluación Sumativa 3}
\newcommand{\portadaSubtitulo}{Modelamiento Predictivo Integrado}
\newcommand{\portadaProyectoUno}{Relación entre condiciones meteorológicas y lluvia al día}
\newcommand{\portadaProyectoDos}{siguiente}
\newcommand{\portadaFecha}{24 de octubre de 2026}
\newcommand{\portadaRepo}{https://github.com/BenjaArayaM/Proyecto-Lluvias-Australia-Grupo7}
\newcommand{\portadaTexto}[3]{%
  \put(297.66,-#1){\makebox[0pt][c]{%
    \fontfamily{ptm}\fontsize{#2bp}{#2bp}\selectfont\color{black}#3}}}
\AddToShipoutPictureFG*{\AtPageUpperLeft{%
  \setlength{\unitlength}{1bp}%
  \put(297.60,-178.80){\makebox[0pt][c]{%
    \includegraphics[width=125.28bp,height=103.92bp]{logo_unab.jpeg}}}%
  \portadaTexto{203.40}{12}{UNIVERSIDAD ANDRÉS BELLO}%
  \portadaTexto{223.32}{11.52}{Magíster en Ciencia de Datos e Inteligencia Artificial}%
  \portadaTexto{251.28}{11.52}{MCDI501}%
  \portadaTexto{265.08}{11.52}{Estadística Computacional para la Toma de Decisiones}%
  \portadaTexto{309.36}{15}{\bfseries\portadaEvaluacion}%
  \portadaTexto{332.64}{12}{\portadaSubtitulo}%
  \portadaTexto{369.72}{18}{\portadaProyectoUno}%
  \portadaTexto{390.48}{18}{\portadaProyectoDos}%
  \portadaTexto{419.64}{11.04}{Informe técnico grupal · Grupo 7 · Ponderación 60 \%}%
  \portadaTexto{451.92}{11.04}{\bfseries Integrantes}%
  \portadaTexto{470.16}{11.04}{Benjamín Araya}%
  \portadaTexto{483.48}{11.04}{Alan Cañete}%
  \portadaTexto{496.80}{11.04}{Patricio Cortez}%
  \portadaTexto{510.00}{11.04}{Jorge Guaico}%
  \portadaTexto{541.32}{11.04}{Dataset: Rain in Australia --- Weather Australia}%
  \portadaTexto{554.64}{11.04}{Docente: Jean Paul Maidana}%
  \portadaTexto{567.84}{11.04}{Fecha: \portadaFecha}%
  \portadaTexto{594.12}{11.04}{\bfseries Repositorio y rama del análisis}%
  \portadaTexto{612.48}{11.04}{%
    \href{\portadaRepo}{\color{blue}\underline{\portadaRepo}}}%
}}
\thispagestyle{empty}
\null
\clearpage


\titulo{1. Resumen ejecutivo e introducción}{Integración de evidencia y decisiones para clasificar RainTomorrow}
El modelo seleccionado mediante validación interna alcanza un ROC-AUC de
\AUCFinal{}, recall de \RecallFinal{}\,\% y precision de \PrecisionFinal{}\,\%
en los 42.658 registros de prueba. El recall aumenta \CambioRecall{} puntos
porcentuales respecto de F2, con un cambio de accuracy de \CambioAccuracy{}
puntos. La mejora implica un intercambio entre detección y falsas alarmas:
todavía quedan \FalsosNegativos{} días lluviosos sin detectar y
\FalsosPositivos{} falsas alertas. El remuestreo identifica incertidumbre
geográfica que el intervalo convencional subestima para algunos coeficientes.

\subtitulo{Problema y alcance}
Se estudia la relación entre condiciones meteorológicas de un día registrado
y \texttt{RainTomorrow}. La aplicación de referencia es una alerta que ayude a
planificar actividades sensibles a lluvia. El archivo \textit{Rain in Australia}
(s.\,f.) contiene 145.460 observaciones, 23 columnas y 49 localidades.
Las 3.267 respuestas ausentes se excluyen; quedan 142.193 registros etiquetados.
Se conserva el archivo original y se documenta su identidad mediante SHA-256.
El análisis es retrospectivo: el dataset no certifica una hora operacional de
emisión ni disponibilidad en tiempo real de todas las mediciones del día.

\begin{center}
\captionof{table}{La evidencia anterior determina decisiones verificables de S3.}
\small
\begin{tabularx}{\linewidth}{@{}lXX@{}}\toprule
Fase & Evidencia propia & Decisión en el cierre\\\midrule
S1 y S2 & Diferencia media de humedad 20,08 pp; IC bootstrap [18,52; 21,62]. & Incorporar \texttt{Humidity3pm} e interpretar su asociación ajustada.\\
S2 & MaxTemp--Temp3pm: $\rho=0,984$, estable bajo remuestreo. & Conservar MaxTemp y excluir Temp3pm de los clasificadores.\\
S1 y S2 & Sunshine ausente en 48,01\,\% del archivo y en todas las filas de 19 localidades. & Comparar eliminación, mediana y regresión; declarar el supuesto de transporte a estaciones sin datos.\\
S2 & Humedad--Rainfall cambia de 0,444 a 0,307 al filtrar en 2 mm. & Conservar extremos plausibles; evaluar log1p(Rainfall) sin eliminar lluvia legítima.\\
F2 & Recall 39,82\,\%, AUC 0,8164 y 5.755 falsos negativos. & Ponderar clases y fijar umbrales con predicciones OOF de entrenamiento.\\\bottomrule
\end{tabularx}
\end{center}
\fuente

\subtitulo{Objetivo integrado}
Comparar estrategias de faltantes y tres procedimientos logísticos, evaluar
estabilidad de coeficientes y traducir la evidencia a recomendaciones.
La partición 70/30 estratificada y la semilla 42 conservan F2: 99.535 registros
de entrenamiento y 42.658 de prueba. Los valores de F2 se reprodujeron antes
del desarrollo. Los antecedentes exploraron el archivo completo y F2 ya observó
la prueba; los resultados describen esta comparación académica y no una
validación externa independiente.
\clearpage

\titulo{2. Manejo de datos faltantes}{Regresión múltiple y elección informada de estrategia}
Se recuperó la auditoría de S1. Entre respuestas conocidas, Sunshine presenta
\FaltantesSunModelo{}\,\% de ausencia. El patrón por localidad cuestiona MCAR;
MAR condicionado a covariables se utiliza como supuesto de trabajo, sin descartar
MNAR (Rubin, 1976). La regresión determinista no elimina esta incertidumbre.

Cada una de las 16 variables numéricas incompletas se modeló mediante OLS con
tres predictores y casos completos de entrenamiento. Se priorizaron pares
validados en S2 y se limitó la redundancia entre predictores. La matriz exportada
de S1 cubre siete variables; las relaciones restantes se completaron y rotularon
como \textit{ampliación S3 de entrenamiento}. La respuesta futura no participa
en ninguna imputación. La validación interna retiene el 20\,\% de donantes.

\begin{center}
\captionof{table}{Reconstrucción de variables numéricas en donantes observados de entrenamiento.}
\TablaImputacion
\end{center}
\fuente

En Sunshine, el RMSE baja de \RMSESolMed{} a \RMSESolReg{} horas
(\MejoraRMSE{}\,\%; $R^2=\RDSol{}$). Se eligió regresión con la regla
predefinida de reducción de RMSE de al menos 5\,\% frente a mediana.
La elección utiliza reconstrucción de entrenamiento y cobertura, no el desempeño
de prueba. Mediana y eliminación se conservan para el análisis de sensibilidad.

\figura{01_imputacion_regresion.png}{4.7cm}{Validación de Sunshine y QQ de residuos: buena reconstrucción no equivale a normalidad ni valida estaciones sin datos.}

Si faltan covariables al aplicar una regresión, se usa su mediana auxiliar de
entrenamiento: ocurrió en \FallbackSol{} de \ImputadosSol{} imputaciones de
Sunshine en train. Los recortes físicos se aplican solo a valores imputados;
los observados permanecen intactos. Shapiro y Breusch--Pagan señalan desviaciones
de normalidad y homocedasticidad; las regresiones se utilizan para predicción
condicional, con cautela inferencial. Los CSV detallan conteos y diagnósticos.
\clearpage

\titulo{3. Construcción y selección de modelos}{Variables informadas por S1/S2 y validación dentro de entrenamiento}
M1 conserva Humidity3pm, Sunshine, Rainfall y MaxTemp de F2. M2 aplica stepwise
bidireccional AIC a ocho candidatos: los antecedentes se complementan en train
con viento, presión, humedad matinal y nubosidad. Log1p(Rainfall) mantiene el
orden de las precipitaciones y modera su asimetría sin descartar extremos.
M3 minimiza BIC entre los 14 subconjuntos de uno a tres predictores de M1.
La restricción hace explícita la solicitud de un modelo reducido: no se presenta
como el BIC mínimo entre todas las especificaciones.

AIC y BIC se calculan sin ponderación, sobre las mismas observaciones, según
sus penalizaciones por complejidad (Akaike, 1974; Schwarz, 1978).
Después se reajusta cada modelo para clasificación con pesos
$w_c=n/(2n_c)$: \PesoYes{} para lluvia y \PesoNo{} para ausencia
(Scikit-learn developers, s.\,f.). Los scores resultantes no están calibrados
automáticamente como probabilidades meteorológicas.

\begin{center}
\captionof{table}{Selección sobre entrenamiento: AIC/BIC no ponderados y métricas fuera de pliegue.}
\TablaSeleccion
\end{center}
\fuente

Los cinco pliegues estratificados reajustan imputación y escalamiento y repiten
la selección de variables dentro del entrenamiento interno. Se escoge el mayor
AUC OOF; diferencias de hasta 0,002 favorecen menos predictores. Cada umbral
maximiza F1 OOF. La especificación y los umbrales quedan registrados antes de
predecir prueba. El procedimiento elegido es \ModeloFinal{}, con
umbral \UmbralFinal{} sobre su score ponderado.

\begin{center}
\captionof{table}{Coeficientes estandarizados y odds ratios de las tres especificaciones.}
\TablaCoeficientesTres
\end{center}
\fuente

En el modelo final, aumentar Humidity3pm en \DEHumedad{} pp se asocia con
multiplicar las odds del ajuste ponderado por \ORHumedad{}, manteniendo
constantes los otros predictores. Para log1p(Rainfall), duplicar $(1+\text{mm})$
se asocia con OR=\ORDoblarRain{}. Un OR no multiplica directamente la probabilidad.
Los signos ajustados pueden diferir de las correlaciones marginales de S1;
no representan efectos causales. El notebook presenta los resúmenes y los
intervalos convencionales de las tres especificaciones.
\clearpage

\titulo{4. Desempeño predictivo}{Discriminación, detección de lluvia y falsas alertas}
\begin{center}
\captionof{table}{Desempeño aparente de entrenamiento y de prueba con umbrales F1 OOF.}
\TablaMetricas
\end{center}
\fuente
\figura{03_modelos_roc_metricas.png}{5.6cm}{ROC y métricas de prueba. El AUC es discriminación y no porcentaje de clasificaciones correctas.}
\figura{04_matrices_confusion.png}{7.7cm}{Matrices para los tres modelos en train y test. Filas reales, columnas predichas; conteos y porcentajes por clase real.}

El modelo final obtiene accuracy \AccuracyFinal{}\,\%, precision
\PrecisionFinal{}\,\%, recall \RecallFinal{}\,\%, F1 \FUnoFinal{} y PR-AUC
\PRAUCFinal{}. Detecta \VerdaderosPositivos{} días de lluvia y omite
\FalsosNegativos{}. El aumento del recall frente a F2 convive con más falsas
alarmas; el uso de alertas requiere valorar ese costo. La mejora cambia
variables, imputación, pesos y umbral simultáneamente, por lo que no se atribuye
a una única intervención. Se publican además métricas con umbral 0,50 como referencia.
\clearpage

\titulo{5. Estabilidad de coeficientes}{Bootstrap de localidades y comparación con intervalos tradicionales}
Se obtuvieron \BootstrapValidas{} remuestras válidas y
\NumeroFallosBootstrap{} fallos, remuestreando las 49 localidades de entrenamiento
con reemplazo. Sus registros disponibles permanecen juntos. En cada remuestra
se reajustan regresiones de imputación, medianas auxiliares, escala y pesos de
clase. Se fijan las listas de predictores: la inferencia es condicional a la
selección y no incorpora su incertidumbre completa. La independencia entre
localidades sigue siendo un supuesto (Cameron y Miller, 2015).

Para acelerar sin modificar el estimador se usan multiplicidades de localidad
y estadísticos OLS por conglomerado. Se verificó la equivalencia contra duplicación
real de filas y un GLM independiente. Las betas se expresan por la misma DE de
referencia, evitando mezclar escalas entre remuestras.

\begin{center}
\captionof{table}{IC 95\,\% de betas y OR. Signo\% es la frecuencia descriptiva de conservar el signo original.}
\TablaBootstrap
\end{center}
\fuente
\figura{05_bootstrap_OR.png}{6.8cm}{Intervalos Wald y percentil del modelo final. OR=1 indica ausencia de asociación ajustada en esta parametrización.}

El intervalo percentil incluye cero para \PredictoresInestables{}; su
correspondiente intervalo de OR incluye uno. La comparación con Wald muestra
por qué significancia convencional y estabilidad geográfica son preguntas
distintas. Humidity3pm conserva una asociación positiva, con OR
\ORHumedad{} e IC bootstrap \ICHumedad{} por \DEHumedad{} pp.
Se calculó también BCa mediante jackknife de 49 localidades y se exportaron
distribuciones y convergencia por prefijos. La frecuencia de signo no es una
probabilidad posterior. El remuestreo no resuelve MNAR ni confusión no medida.
\clearpage

\titulo{6. Diagnóstico y alcance inferencial}{Multicolinealidad, forma funcional, residuos e influencia}
\begin{minipage}[t]{0.47\linewidth}
\captionof{table}{VIF en entrenamiento preparado.}
\TablaVIF
\end{minipage}\hfill
\begin{minipage}[t]{0.50\linewidth}
\captionof{table}{Box--Tidwell exploratorio y cambio relativo de deviance.}
\TablaLinealidad
\end{minipage}
\par\medskip
El VIF máximo es \MaxVIF{}, sin evidencia de redundancia extrema como la del
par MaxTemp--Temp3pm descartado. Las extensiones no lineales detectan desviaciones
para algunas variables; sus reducciones relativas de deviance ayudan a separar
significancia con gran muestra de importancia práctica. Los valores p son
exploratorios bajo el ajuste ponderado y no corrigen la dependencia espacial.
Estos diagnósticos no permiten afirmar que la forma lineal del logit sea exacta.

\figura{07_diagnostico.png}{6.2cm}{Residuos y calibración OOF. Un AUC favorable puede coexistir con scores que no coinciden con frecuencias observadas.}

La regresión logística no exige residuos normales ni homocedásticos. Se
examinaron residuos de deviance y Pearson, leverage y Cook. El umbral descriptivo
$4/n$ marca \CookMarcados{} registros; no se eliminaron automáticamente.
La magnitud de las señales y el contexto meteorológico deben acompañar cualquier
decisión de depuración. Conservar extremos plausibles responde a la sensibilidad
del filtro de S2, que cambiaba la muestra al cortar precipitaciones en 2 mm.

\subtitulo{Influencia por localidad}
El jackknife identifica \LocalidadesInfluyentes{} como las tres localidades
con mayor cambio máximo de beta en escala común. Esto indica influencia sobre
el ajuste, sin certificar errores de medición. Se publican las 49 omisiones y
las 20 observaciones con mayor Cook para revisión trazable.

\subtitulo{Interpretación de probabilidades}
La ponderación de clases modifica el objetivo del ajuste. La curva de calibración
OOF evidencia la distancia entre score y frecuencia real de lluvia. En este trabajo
el score sirve para ordenar y clasificar con un umbral validado internamente;
su uso como probabilidad operacional requiere calibración y validación adecuadas
al horizonte y al contexto de decisión.
\clearpage

\titulo{7. Impacto de la imputación}{Misma especificación, umbral y escala de interpretación}
\begin{center}
\captionof{table}{Comparación sobre prueba propia y sobre los mismos casos comunes.}
\TablaImpacto
\end{center}
\fuente
La eliminación cubre \CoberturaEliminacion{}\,\% de prueba y
\LocEliminacion{} localidades. Un AUC superior en esa muestra no demuestra
superioridad sobre estrategias que cubren las 49 localidades. Los casos comunes
permiten comparar procedimientos sobre idénticos registros; aun así, describen
una subpoblación con mayor completitud de medición.

\figura{08_impacto_imputacion.png}{5.1cm}{Cobertura y AUC sobre casos comunes. El tamaño y la composición de la población importan para interpretar la diferencia.}

\begin{center}
\captionof{table}{Betas e IC Wald del modelo final por la misma DE de referencia.}
\TablaCoefImpacto
\end{center}
\fuente

La imputación por mediana concentra Sunshine y reduce su variación. La regresión
conserva variación asociada a covariables, aunque produce encogimiento y depende
de supuestos de transporte. La estrategia principal se eligió por reconstrucción
de entrenamiento y conservación de cobertura. La comparación predictiva constituye
un análisis de sensibilidad de esa decisión, manteniendo predictores y umbral.
La ausencia estructurada y la incertidumbre sobre valores no observados impiden
interpretar cualquier imputación como recuperación de la verdad meteorológica.
\clearpage

\titulo{8. Simulación para una decisión de planificación}{Ampliación S3 que incorpora la retroalimentación de S2}
La simulación de S2 describía una media. En el cierre se formula una decisión
explícita: cuánta capacidad para atender días lluviosos se necesita en un escenario
hipotético de 30 días, aceptando 5\,\% de riesgo de excedencia. El riesgo es un
supuesto de planificación, no un requisito operacional observado.

Se estiman frecuencias solo en entrenamiento con Humidity3pm observado:
\ProbMCAlta{}\,\% cuando la humedad es al menos 70\,\% y
\ProbMCBaja{}\,\% cuando es menor. Bajo días iid y parámetros fijos,
$D\sim\mathrm{Binomial}(30,\widehat p)$. Se generan 10.000 escenarios por grupo,
se contrasta la media con $30\widehat p$ y se toma el cuantil discreto superior
del 95\,\% como capacidad de referencia.

\begin{center}
\captionof{table}{Calibración de entrenamiento, resultados Monte Carlo y capacidad de referencia.}
\TablaMonteCarlo
\end{center}
\fuente
\figura{09_montecarlo_decision.png}{6.2cm}{Distribución del número de días lluviosos y capacidad que cubre al menos el 95\,\% de los escenarios simulados.}

Las medias son \MediaMCAlta{} y \MediaMCBaja{} días; las capacidades resultan
\CapacidadAlta{} y \CapacidadBaja{} días, respectivamente. El resultado informa
una reserva de recursos bajo condiciones estipuladas y hace explícita la tolerancia
al riesgo. Las tablas registran convergencia, MCSE y probabilidad analítica de
exceder 15 días. Una ausencia de excedencias en 10.000 simulaciones no implica
probabilidad nula: se informa también la cota binomial de Monte Carlo.

\subtitulo{Discusión integrada}
S1 caracteriza asociaciones y faltantes; S2 evalúa su estabilidad y muestra
sensibilidad al rango de lluvia; F2 establece una referencia predictiva;
S3 compara procedimientos, detecta coeficientes inestables y explicita decisiones.
La simulación no pronostica 30 fechas futuras: ignora dependencia diaria,
estacionalidad y parte de la incertidumbre de $p$. La partición aleatoria exigida
evalúa la muestra disponible. La generalización temporal y geográfica queda
como limitación, conforme al alcance indicado por el docente.
\clearpage

\titulo{9. Conclusiones y reproducción}{Recomendaciones sustentadas por el proyecto completo}
\begin{itemize}
\item \textbf{Detección y errores.} El procedimiento \ModeloFinal{} alcanza AUC
\AUCFinal{} y recall \RecallFinal{}\,\%. Su uso debe valorar conjuntamente
\FalsosNegativos{} omisiones y \FalsosPositivos{} falsas alertas; accuracy
aislada no sustenta una decisión de alerta.
\item \textbf{Cobertura y faltantes.} Se recomienda conservar la comparación con
mediana y comunicar la dependencia de Sunshine respecto de estación. La eliminación
reduce cobertura a \CoberturaEliminacion{}\,\%; la imputación no sustituye medición.
\item \textbf{Incertidumbre.} La estabilidad de Humidity3pm contrasta con
\PredictoresInestables{}. El intervalo por localidad complementa la inferencia
convencional y mantiene visibles sus supuestos.
\item \textbf{Aplicación.} Antes de operar se requiere definir horizonte,
disponibilidad de variables, costos de errores y calibración. La capacidad
\CapacidadAlta{}/\CapacidadBaja{} corresponde únicamente al escenario simulado.
\end{itemize}

\subtitulo{Trazabilidad y reproducción}
El notebook maestro declara semillas 42 y 50101--50104, cinco pliegues y
2.000 remuestras válidas. El código, tablas, figuras y JSON permiten reconstruir
los resultados; las evidencias separan Work y Windows. El SHA-256 del original es
{\fontsize{8.6}{10}\selectfont\texttt{573fd715cd69fcacc4df32024d823b450ae3edaae7e8ff2eeb623adbed424014}}.
Las instrucciones de instalación y la auditoría de los seis criterios acompañan
el repositorio. La participación real en el video y la entrega en Canvas requieren
sus evidencias específicas.

\subtitulo{Referencias}
{\fontsize{9}{11}\selectfont
\setlength{\parskip}{5pt}
\hangindent=1.2em\hangafter=1 Akaike, H. (1974). A new look at the statistical model identification.
\textit{IEEE Transactions on Automatic Control, 19}(6), 716--723.
\url{https://doi.org/10.1109/TAC.1974.1100705}

\hangindent=1.2em\hangafter=1 Cameron, A. C., \& Miller, D. L. (2015). A practitioner's guide to cluster-robust inference.
\textit{Journal of Human Resources, 50}(2), 317--372.
\url{https://doi.org/10.3368/jhr.50.2.317}

\hangindent=1.2em\hangafter=1 \textit{Rain in Australia}. (s.\,f.). Kaggle.
\url{https://www.kaggle.com/datasets/jsphyg/weather-dataset-rattle-package}

\hangindent=1.2em\hangafter=1 Rubin, D. B. (1976). Inference and missing data.
\textit{Biometrika, 63}(3), 581--592.
\url{https://doi.org/10.1093/biomet/63.3.581}

\hangindent=1.2em\hangafter=1 Schwarz, G. (1978). Estimating the dimension of a model.
\textit{The Annals of Statistics, 6}(2), 461--464.
\url{https://doi.org/10.1214/aos/1176344136}

\hangindent=1.2em\hangafter=1 Scikit-learn developers. (s.\,f.). \textit{compute\_class\_weight}.
\url{https://scikit-learn.org/stable/modules/generated/sklearn.utils.class_weight.compute_class_weight.html}
}
\end{document}
